% 13_resources.tex -- Part 13, computational cost. Sources: research_monograph/appendices/MH_resources.tex
% (tab:mono-res-atlas, launch budgets, recorded totals, tab:mono-res-dft, memory, infrastructure, energy estimate);
% evidence/derived/atlas_usage.json (launches, wall time, non-zero exits, longest launch);
% research_record/05_DECISION_LOG.md (D-02 budget 52; D-03 budget 57 -> 71; D-13 200 more; D-58 C10 17/200, cap not
% raised); appendices/MG_corrections.tex T1 (C10 runs from 0061, 17 of 200 by 2026-10-14);
% chapters/13_dft_rev.tex sec:dft-resources (totals, laptop, 30-100 GB) and sec:dft-gpu line 1421 (GPU benchmark:
% the same 96-atom SCF, 9 min 16 s on one H100 NVL against 1 h 16 min on 16 CPU cores, factor 8.2; NOTE: the row
% pairing in tab:dft-timing, 1 h 48 min vs 9 min, mixes the 35 A and 15 A vacuum cells and is not used here);
% chapters/M04b_dft_numerics.tex sec:mono-dft-cost (N^3 scaling, 328 -> 680 bands, 3.6 GB, 82 GB for the
% single-process GPU relaxation); qe_workflow/cern_htcondor/results/resources.json (pool snapshot 2026-10-06) and the
% PWSCF WALL lines and GPU model lines of the scf/relax/bands outputs in qe_workflow/cern_htcondor/results/ (wall
% times and hardware re-checked 2026-10-15; every GPU row ran on an H100 NVL).
% Energy: ESTIMATE with assumed power figures (MH), not an audit. No sealed run (0075-0077) and no amorphous
% electronic quantity appears. Tables and one TikZ diagram only; no figure files.
% Note: frame bodies start with a command, never with a brace group (beamer would read a subtitle).
\providecommand{\thirteendense}{\setbeamerfont{itemize/enumerate body}{size=\footnotesize}%
  \setbeamerfont{itemize/enumerate subbody}{size=\footnotesize}}
\providecommand{\thirteenTop}{\vspace{-4mm}}
\providecommand{\thirteenAtlas}{\footnote[frame]{Silvaco, \citetitle{AtlasManual}}}
\providecommand{\thirteenNpw}{N_{\mathrm{pw}}}
\providecommand{\thirteenNb}{N_{\mathrm{b}}}
\providecommand{\thirteenNk}{N_{k}}
\providecommand{\thirteenHead}[1]{{\footnotesize\bfseries #1\par}\vspace{0.8mm}}
\providecommand{\thirteenX}{\,$\times$\,}
\providecommand{\thirteenx}{$\times$}

\partslide{Part 13}{Computational cost}

% ---------------------------------------------------------------------------------------------------------------------
\begin{frame}[t]{Device simulation is cheap; launch budgets limit it}
  \thirteenTop\thirteendense
  \begin{columns}[T,onlytextwidth]
    \begin{column}{0.58\textwidth}
      \thirteenHead{ATLAS launches in the public record\thirteenAtlas}
      {\footnotesize\setlength{\tabcolsep}{4.5pt}
      \begin{tabular}{@{}lrrr@{}}
        \toprule
        Study & Launches & \shortstack[r]{Wall\\time (h)} & \shortstack[r]{Exit\\$\neq 0$} \\
        \midrule
        Model V1 study      & 57 & 7.8 & 2 \\
        Model V2            &  3 & 0.5 & 0 \\
        Campaign C10        & 14 & 2.9 & 1 \\
        \midrule
        All public          & 74 & 11.1 & 3 \\
        \bottomrule
      \end{tabular}\par}
      \src{Silvaco ATLAS 5.28, 32-bit, local licence. Runs 0001--0057, 0058--0060, 0061--0074; wall times from
        \texttt{results/RUN\_INDEX.csv}; exit $\neq 0$: non-zero exit code; sealed runs 0075--0077 not included.
        Budgets: decision log D-02, D-03, D-13, D-58.}
    \end{column}
    \begin{column}{0.39\textwidth}
      \vspace{4.2mm}
      \begin{itemize}
        \item Launches run one at a time; the longest took 1.0\,h \lab{measured}. Each is a separately reviewed
          decision.
        \item Campaign C10 had used 17 launches by 14 October: the 14 in the table and the 3 sealed runs.
      \end{itemize}
    \end{column}
  \end{columns}
  \thirteenHead{ATLAS launch budget in the decision log (in order of date, not to scale)}
  \begin{tikzpicture}[x=1cm,y=1cm,font=\footnotesize]
    \foreach \a/\b in {0.25/3.15, 3.15/6.05, 6.05/10.0}
      \draw[rulegrey, line width=0.8pt, -{Stealth[length=1.6mm]}] (\a+0.15,0) -- (\b-0.17,0);
    \foreach \x/\d/\txt in {0.25/25 Sep/{budget: 52 launches},
        3.15/5 Oct/{budget 57 $\to$ 71},
        6.05/10 Oct/{200 more launches for C10}} {
      \fill (\x,0) circle (0.065);
      \node[anchor=south west, inner sep=0] at (\x-0.1,0.2) {\textbf{\d\ 2026}};
      \node[anchor=north west, inner sep=0] at (\x-0.1,-0.22) {\txt};
    }
    \fill[ntumaroon] (10.0,0) circle (0.085);
    \node[anchor=south west, inner sep=0] at (9.9,0.2) {\textbf{14 Oct 2026}};
    \node[anchor=north west, inner sep=0] at (9.9,-0.22)
      {\textcolor{ntumaroon}{17 of 200 used}; cap not raised};
  \end{tikzpicture}
  \takeaway{Device simulation costs hours; the scarce resource is the reviewed launch.}
\end{frame}

% ---------------------------------------------------------------------------------------------------------------------
\begin{frame}[t]{DFT to 4 October: about 8\,300 core-hours, 85--90 GPU-hours}
  \thirteenTop\thirteendense
  \begin{columns}[T,onlytextwidth]
    \begin{column}{0.545\textwidth}
      \thirteenHead{Quantum ESPRESSO\srefs{Giannozzi2009,Giannozzi2017} runs: measured wall times}
      {\scriptsize\setlength{\tabcolsep}{3pt}
      \begin{tabular}{@{}llr@{}}
        \toprule
        Step & Resources & Wall time \\
        \midrule
        80-atom W cell, relaxation                 & 32 CPU cores   & 23\,h 48\,min \\
        80-atom W cell, SCF, 4\thirteenx4\thirteenx4 $k$ & 32 CPU cores & 3\,h 03\,min \\
        1\,nm W film (96 atoms), SCF               & 16 CPU cores   &  2\,h 55\,min \\
        1\,nm film, bands                          & 16 CPU cores   &  4\,h 38\,min \\
        2\,nm W film (176 atoms), relax, 28 BFGS   & one H100 NVL   & 16\,h 49\,min \\
        2\,nm W film, final SCF                    & 33 CPU cores   &  8\,h 05\,min \\
        2\,nm film, bands                          & one H100 NVL   &  8\,h 22\,min \\
        0.695\,nm film, relaxation                 & one H100 NVL   &  9\,h 58\,min \\
        HSE06, 1\,nm film, SCF, 4 $k$ points       & one H100 NVL   &  4\,h 00\,min \\
        HSE06, bulk, SCF, 3\thirteenx3\thirteenx3 $q$ mesh & one H100 NVL & 4\,h 15\,min \\
        \bottomrule
      \end{tabular}\par}
      \src{Representative steps; PWSCF wall times in the output files. The 2\,nm W film and the 0.695\,nm film belong to the later
        jobs and are not in the totals.}
    \end{column}
    \begin{column}{0.425\textwidth}
      \thirteenHead{Recorded usage up to 4 October 2026}
      \begin{itemize}
        \item CERN: about 8\,300 weighted CPU core-hours and 85--90 GPU-hours; laptop: about 190 core-hours
          \lab{derived}.
        \item About 3\,000 of the core-hours and 50 GPU-hours were lost to failures, since fixed \lab{derived}.
        \item A laptop would need about six weeks for the useful work \lab{inferred} and could not run the 2\,nm
          and hybrid films (30--100\,GB).
        \item Energy estimate (not an audit) at 10\,W per core and 0.4\,kW per GPU \lab{assumed}:
      \end{itemize}
      \vspace{0.3mm}
      {\footnotesize\centering 8\,300\,h\thirteenX10\,W + 90\,h\thirteenX0.4\,kW $\approx$ 120\,kWh\par}
    \end{column}
  \end{columns}
  \takeaway{Estimated energy to 4 October: about 120\,kWh; later jobs are not yet totalled.}
\end{frame}

% ---------------------------------------------------------------------------------------------------------------------
\begin{frame}[t]{At 66--82\,GB, memory decided where the 176-atom film ran}
  \thirteenTop\thirteendense
  \begin{columns}[T,onlytextwidth]
    \begin{column}{0.53\textwidth}
      \thirteenHead{Plane-wave cost for a cell of $N$ atoms}
      {\centering$t_{\mathrm{it}} \propto \thirteenNb^{2}\,\thirteenNpw \propto N^{3}, \qquad
         M_{\psi} \propto \thirteenNb\,\thirteenNpw\,\thirteenNk \propto N^{2}\thirteenNk$\par}
      \vspace{0.8mm}
      {\scriptsize\color{inkgrey} $t_{\mathrm{it}}$: time per SCF iteration; $M_{\psi}$: wave-function memory;
        $\thirteenNb$, $\thirteenNpw$, $\thirteenNk$: bands, plane waves, $k$ points. Linear-scaling
        methods\sref{Goedecker1999} not used.\par}
      \vspace{1mm}
      \begin{itemize}
        \item 1\,nm $\to$ 2\,nm film: bands 328 $\to$ 680, plane waves +29\,\%; wave functions \thirteenx2.7,
          $\thirteenNb^{2}\thirteenNpw$ work \thirteenx5.5 \lab{derived}.
        \item 2\,nm film: 680\thirteenX327\,585\thirteenX16\,B = 3.6\,GB per $k$ point \lab{derived};
          \texttt{pw.x} estimates 82\,GB for the single-process GPU relaxation \lab{computed}; the final SCF needed
          66\,GB \lab{measured}. Hence a large-memory GPU or a distributed CPU run.
        \item HSE06, 1\,nm film: the default exact-exchange cutoff ran out of memory; 160\,Ry was used
          \lab{observed}.
      \end{itemize}
    \end{column}
    \begin{column}{0.44\textwidth}
      \thirteenHead{CERN GPU pool, 6 Oct 2026 \lab{observed}}
      {\footnotesize
      \begin{tabular}{@{}lr@{}}
        \toprule
        GPU model & GPUs \\
        \midrule
        H100 NVL & 144 \\
        H200 & 96 \\
        A100-PCIE-40GB & 72 \\
        Older (V100, V100S, T4) & 93 \\
        \midrule
        Total, on 79 nodes & 405 \\
        \bottomrule
      \end{tabular}\par}
      \src{HTCondor snapshot (\texttt{resources.json}); jobs from lxplus, data on EOS.}
      \vspace{1mm}
      \begin{itemize}
        \item A pure 1\,nm film SCF benchmark (96 atoms): 1\,h 16\,min on 16 CPU cores, 9\,min on one H100 NVL (\thirteenx8.2)
          \lab{measured}.
        \item About 0.8 GPU starts per hour \lab{observed}; two-GPU requests of the amorphous campaign matched no
          slot.
      \end{itemize}
    \end{column}
  \end{columns}
  \takeaway{For the large cells, memory and access to large GPUs, not core-hours, were the limits.}
\end{frame}
